(a) Every even-length term ends in $67$, hence is three modulo four and cannot be square. The three-digit term is $676=26^2$. Every odd-length term of length at least five ends in $7676$, which is twelve modulo sixteen, not a square residue. Thus $\boxed{676}$ is the only square. (b) Put $c=10a+b$. The $2k$-digit term is $\dfrac{c(100^k-1)}{99}$, and the $2k+1$-digit term is $\dfrac{10c\,100^k-(a+10b)}{99}$, with $k\ge1$ in the even case and $k\ge0$ in the odd case. Consequently the necessary and sufficient tests are respectively
\[99s^2+c=c100^k,\qquad99s^2+a+10b=10c100^k\]
for an integer $s\ge0$. These follow by summing the place-value geometric series and are exact for every specified pair and length. They also give useful last-digit modular obstructions, as in part (a); they are not an asserted finite classification of all digit pairs over all lengths.
